\capitulo{4}{Técnicas y herramientas}

Esta parte de la memoria tiene como objetivo presentar las técnicas metodológicas y las herramientas de desarrollo que se han utilizado para llevar a cabo el proyecto. Si se han estudiado diferentes alternativas de metodologías, herramientas, bibliotecas se puede hacer un resumen de los aspectos más destacados de cada alternativa, incluyendo comparativas entre las distintas opciones y una justificación de las elecciones realizadas. 
No se pretende que este apartado se convierta en un capítulo de un libro dedicado a cada una de las alternativas, sino comentar los aspectos más destacados de cada opción, con un repaso somero a los fundamentos esenciales y referencias bibliográficas para que el lector pueda ampliar su conocimiento sobre el tema.


# Conceptos Técnicos:

Entorno de desarrollo: Implementación del laboratorio virtual mediante Unity, integrando prefabs (bloques y piezas lego) y simulaciones en tiempo real
Modelado 3D: Diseño y exportación de componentes robóticos con Blender para integrarlos en Unity, teniendo en cuenta la optimización de texturas y mallas.
Programación orientada a bloques: Recreación de un motor de programación visual basado en Blockly/Scratch, integrando funcionalidades específicas para la simulación y control de robots.
Control de versiones: Integración de Plastic SCM para gestión de cambios y sincronización con repositorios en Unity Cloud y GitHub, asegurando la trazabilidad del desarrollo.
>Bases de datos en tiempo real:</strong> Uso de Firebase para gestionar y sincronizar datos entre usuarios, permitiendo almacenar información sobre proyectos, configuraciones y resultados en tiempo real.
Diseño y modelado UML: Creación de diagramas UML para representar las interacciones del sistema, el flujo de datos y los componentes clave del software.
Compatibilidad con dispositivos VR: Diseño del entorno 3D para garantizar una experiencia fluida en gafas de RV como Oculus, considerando el rendimiento y la optimización para hardware específico.
Validación y pruebas de software: Metodologías para garantizar la calidad del software, incluyendo pruebas unitarias para validar módulos individuales, pruebas de integración para garantizar la interoperabilidad de los componentes, y pruebas de usabilidad para evaluar la experiencia del usuario objetivo.

Aplicación de conceptos relacionados con metodologías ágiles.


